% related.tex
% Four paragraphs, tightened in the language-audit pass (2026-07-14;
% see project_language_audit_2026-07-14 memory): (1) parametric /
% stylized-facts generators; (2) latent-state and deep generative
% approaches with the hybrid HMM contribution; (3) SIM composition
% (full-return vs residual-fit) with the variance-correction
% contribution; (4) validation including the diversification-return
% concern. Cut ~200 words of historical survey and repeated
% "captures X but fails Y" phrasing; removed "Methods describes"/
% "developed in Methods" section self-references (recast with "we").

Approaches to reproducing the stylized facts of equity markets
divide into bottom-up agent-based models~\cite{luxMarchesi1999,
contBouchaud2000, farmerJoshi2002, lebaron2006, arthur2021,
farmer2025}, which derive return dynamics from simulated
interactions of heterogeneous agents but are demanding to calibrate
and expensive to simulate at scale, and top-down models, which fit
statistical structures directly to observed growth-rate series.
Top-down generators trace to Bachelier's diffusion model, later
formalized as geometric Brownian motion and used by
Black-Scholes-Merton for option
pricing~\cite{bachelier1900, osborne1959, samuelson1965,
blackScholes1973}; these diffusion baselines are tractable but fail
on the empirical regularities that drive risk and pricing. Mandelbrot~\cite{mandelbrot1963}
showed returns have heavier tails than the Gaussian admits;
Fama~\cite{famaEMH1970} reviewed evidence that raw returns have
near-zero autocorrelation; and Schwert~\cite{schwert1989} documented
persistence in monthly absolute residuals, later codified by
Cont~\cite{cont2001} as volatility clustering. The autoregressive
conditional heteroskedasticity (ARCH) and generalized ARCH (GARCH)
models of Engle and Bollerslev~\cite{engle1982, bollerslev1986}
reproduced the clustering but left marginals light-tailed unless the
noise term was itself heavy-tailed, and neither models discrete
regimes or sudden jumps. Stochastic-volatility
models~\cite{steinStein1991, heston1993} make volatility a latent
process to capture clustering, but estimating that process requires
filtering or Markov chain Monte Carlo (MCMC)
methods~\cite{kimShephardChib1998}; jump-diffusion
models~\cite{merton1976, kou2002} inject Poisson events into the
price path, producing heavy tails by construction, but volatility
does not persist after the jump.

Hidden Markov models embed regime and jump behavior directly in a
latent discrete chain, following Hamilton's~\cite{hamilton1989}
two-state regime-switching formulation, later applied to volatility
forecasting~\cite{rossiGallo2006} and equity
trading~\cite{nguyen2018}. Ryden et
al.~\cite{rydenTerasvirtaAsbrink1998} found that Gaussian-emission
HMMs matched heavy tails and the absence of return autocorrelation
but failed on volatility clustering; Bulla and
Bulla~\cite{bullaBulla2006} recovered the clustering with hidden
semi-Markov models (HSMMs) that replaced the geometric dwell-time
distribution of every state with a parametric (negative-binomial)
alternative. Both rely on the Baum-Welch
expectation-maximization (EM) algorithm~\cite{baumPetrieSoulesWeiss1970},
which is iterative and sensitive to initialization. Deep generative
alternatives learn the marginal directly but routinely miss the
volatility-clustering signature: Takahashi et
al.~\cite{takahashiEtAl2019} found that generative adversarial
network (GAN) architectures produced return sequences that resembled
financial data yet failed to reproduce the absolute-return
autocorrelation function, and Kwon and Lee~\cite{kwonLee2024}
confirmed the same on TimeGAN~\cite{yoonEtAl2019timegan}. Recent
surveys~\cite{assefaEtAl2020, jordonEtAl2022, stengerEtAl2024} agree
that stylized-facts reproduction remains the primary quality
criterion. This model also uses a latent-state chain but follows a
different design principle than the HSMM: rather than replace every
state's dwell-time distribution, it keeps the empirically counted
transition matrix for the central states and adds an event-driven
override that forces the chain into the tail states for a
Poisson-distributed duration on a subset of time steps.

Multi-asset composition adds cross-asset dependence on top of a
univariate generator. Mixture hidden Markov
models~\cite{diasVermuntRamos2010} share latent state across assets
but are costly to estimate in high dimensions; copula
approaches~\cite{embrechtsMcNeilStraumann2002,
cherubiniLucianoVecchiato2004, aasCzadoFrigessiBakken2009} separate
marginals from cross-asset dependence but require choosing a copula
family and, in their standard static form, treat successive time
slices as independent. Factor models instead reduce a firm's growth
rate to one or more shared factors~\cite{rossAPT1976,
famaFrench1993}; the single-index model (SIM) of
Sharpe~\cite{sharpe1963} is the simplest, decomposing each asset's
growth rate into a market-loaded term plus an idiosyncratic
residual, a linear, easy-to-interpret structure that pairs naturally
with a single-factor market generator. Within SIM, the choice of
residual specification trades off the per-asset generator's
distributional structure against the composed variance; we derive a
closed-form variance correction that recovers both.

Validation of synthetic growth-rate paths combines two families of
tests: goodness-of-fit tests for distributional
fidelity~\cite{kolmogorov1933, smirnov1948, andersonDarling1952} and
autocorrelation diagnostics for temporal structure~\cite{cont2001}.
These tests are necessary but not sufficient: a synthetic dataset
can pass them yet still mislead when fed into the task it was
generated for. Booth and Fama~\cite{boothFama1992} described one
such failure: a daily-rebalanced allocator on paths whose residuals
are drawn independently across assets extracts a structural
\emph{diversification return} that is muted in real markets, where
common sector and style factors correlate residuals across assets.
We adopt the classical distributional metrics together with a
Value-at-Risk coverage check scored by Kupiec's
unconditional-coverage test, and address this diversification-return
bias with a location-shift correction.
